% AUTORIDAD DE FUENTE: Exposé IX, §§7--7.1.6, líneas 1845--1887;
% índices de escaneo 410--413, folios 399--402. Las tres lecturas
% materiales D/C del §7.1.3 se conservan reversiblemente mediante \sisi.

\medskip
\noindent
7. \underline{La extensión de Raynaud $G^\natural$ sobre $S$ asociada
al esquema abeliano $A_K$ con reducción semiestable. Dualidad de los
esquemas abelianos $B,B'$ sobre $S$.}

\medskip
\noindent
7.1. Sea $S$ un esquema noetheriano, espectro de un anillo $V$ separado
y completo para la topología definida por un ideal $\underline m$ de
$V$. Sea $A$ un esquema en grupos liso sobre $S$, sea $A^o$ su
componente neutra, y supongamos que $A_o/A_o^o$ es finito (y por tanto
finito étale) sobre $S_o$, y que $A_o^o$ es una extensión de un esquema
abeliano por un toro. Como de costumbre, hemos puesto
$A_n=A\times_S S_n$, donde
$S_n=\operatorname{Spec}(V/\underline m^{n+1})$. Nos interesará
principalmente el caso en que $S$ es un rasgo, la fibra genérica $A_K$
es un esquema abeliano con reducción semiestable sobre $S$, y $A$ es el
modelo de Néron de $A_K$.

Sea $T_o$ el toro maximal de $A_o$, de modo que
$B_o=A_o^o/T_o$ sea un esquema abeliano. Se aplica entonces la
construcción de 5.1 y da un subtoro formal $\widehat T$ del completado
formal $\widehat A$.

%%%%%%%%%%  scan idx 411  |  folio 400  %%%%%%%%%%

Para todo $n$, el cociente
\begin{equation}\tag{7.1.1}
 B_n=A_n^o/T_n
\end{equation}
es representable (SGA~3 $\mathrm{VI}_A$~3.2); más precisamente, es un
esquema plano de tipo finito sobre $S_n$ cuya imagen inversa sobre $S_o$
es el esquema abeliano $B_o$. Por consiguiente, $B_n$ es un
\underline{esquema abeliano} sobre $S_n$. Al variar $n$, los $B_n$ se
obtienen unos de otros por cambio de base y definen, por tanto, un
\underline{esquema abeliano formal}
\begin{equation}\tag{7.1.2}
 \widehat B=(B_n)_{n\geqslant0}\,.
\end{equation}
Daremos condiciones que implican que este esquema abeliano formal es
\underline{algebraizable} (EGA~III~5.4.1), más precisamente, que es el
completado formal de un esquema abeliano proyectivo sobre $S$, que
denotaremos entonces por $B$. Según loc. cit., $B$ está determinado
salvo isomorfismo único (como $S$-esquema; la estructura de grupo de
$\widehat B$ proviene entonces de manera única de una estructura de
grupo sobre $B$, de nuevo por loc. cit.). Este esquema abeliano $B$ se
llamará entonces la \underline{parte abeliana} del esquema $A$, y se
denotará a veces por $A^{\mathrm{ab}}$.

%%%%%%%%%%  scan idx 412  |  folio 401  %%%%%%%%%%

\medskip
\noindent
7.1.3. Por EGA~III~5.4.5, la proyectividad de $\widehat B$ significa
que se puede encontrar un haz invertible $\widehat{\mathbb M}$ sobre
$\widehat B$, es decir, una sucesión de haces invertibles $\mathbb M_n$
sobre los $B_n$ que se pegan al variar $n$, y que sea relativamente
amplio sobre $S$, es decir, tal que $\mathbb M_o$ sea amplio. La idea
natural sería partir de un haz invertible amplio
$\sisi{\mathbb M}{\mathbb L}$ sobre $\sisi{G}{A}$ y demostrar que, para
todo $n$, su restricción $\mathbb L_n$ a $\sisi{G_n}{A_n}$ es la imagen
inversa de un haz invertible $\mathbb M_n$ sobre $B_n$. Ignoro si tales
$\mathbb M_n$ existen necesariamente (y lo dudo mucho). Seguiremos un
método ligeramente distinto, construyendo una \underline{polarización}
de $\widehat B$, es decir, un sistema proyectivo de polarizaciones
$\rho_n$ de los $B_n$ que se pegan. Se sabe [5] que los dobles de estas
polarizaciones definen canónicamente haces invertibles (simétricos)
$\mathbb M_n$ sobre los $B_n$. Por su naturaleza canónica, estos haces
se pegan necesariamente y proporcionan el $\widehat{\mathbb M}$ buscado.

Interpretaremos un elemento de $\operatorname{NS}(B_n/S_n)$ como una
correspondencia divisorial simétrica sobre $(B_n,B_n)$ [5], es decir,
como una biextensión de $(B_n,B_n)$ por $\mathbb G_{mS_n}$. Por VIII~3.5,
estas biextensiones corresponden biyectivamente a biextensiones de
$(A_n,A_n)$ por $\mathbb G_{mS_n}$. Para obtener un sistema coherente de
tales biextensiones $W_n$, basta partir de
\begin{equation}\tag{7.1.4}
 W,\quad\text{una biextensión simétrica de $(A,A)$ por $\mathbb G_{mS}$},
\end{equation}
y tomar los $W_n$ obtenidos a partir de $W$ mediante los cambios de base
$S_n\longrightarrow S$. Debemos suponer además que el elemento de
$\operatorname{NS}(B_o)$ definido por $W_o$ es una polarización de
$B_o$; obtenemos así una polarización de $\widehat B$.

%%%%%%%%%%  scan idx 413  |  folio 402  %%%%%%%%%%

\medskip
\noindent
\underline{Proposición 7.1.5.} \underline{Con las notaciones precedentes
para $S,A,\widehat B$, supongamos que $S$ es normal. Entonces
$\widehat B$ es algebraizable, y el esquema abeliano $B$ sobre $S$ del
que proviene es proyectivo sobre $S$.}

Podemos suponer evidentemente que $A=A^o$ y que $S$ es integral. Por
Raynaud [23, XI~1.13], la hipótesis de que $S$ sea normal implica que
$A$ es cuasiproyectivo sobre $S$. Elijamos un haz invertible
relativamente amplio $\mathbb L$ sobre $A$, y formemos el haz invertible
birigidificado
\[
 \delta(\mathbb L)=
 \pi^*(\mathbb L)\operatorname{pr}_1^*(\mathbb L)^{-1}
 \operatorname{pr}_2^*(\mathbb L)^{-1}
\]
sobre $A\times_S A$, como en VIII~4.12. Puesto que las fibras de $A_o$
tienen rango unipotente cero, lo mismo vale para las fibras de $A$ (esto
se deduce fácilmente de 3.1). Por tanto, según VIII~7.5,
$\delta(\mathbb L)$ define una biextensión $W$ de $(A,A)$, evidentemente
simétrica, que por VIII~4.13 (es decir, esencialmente por Raynaud
[23, XI~1.11]) satisface la condición de amplitud requerida en los
puntos de $S_o$.

\medskip
\noindent
\underline{Observación 7.1.6.} Sin suponer que $S$ sea normal, se puede
demostrar que $\widehat B$ siempre es algebraizable, reduciéndose al
caso normal como en 13.5.3 a).
